\begin{description}
	\item[dataブロック] 群の数\verb|Lv|，各群のサンプルサイズ\verb|N|,各データ点\verb|X|を入力します。ここで\verb|X|が\textbf{二元配列}になっていることに注目してください。\verb|X[1,1]|は第一群の第1番目のデータ，\verb|X[1,2]|は第一群の第2番目のデータ，\verb|X[2,1]|は第二群の第1番目のデータ，同様に\verb|X[i,j]|は第i群の第j番目のデータを表すことになります。この形式は二元配列(2次元の変数セット)になっていると言います\footnote{行列でもいいじゃないか，と思われるかもしれません。それでも結構です。Stanは行列の時，\verb|real|ではなく\verb|matrix|で宣言することができます。しかしこのように\verb|real|型にしておいて，後ろのカッコで配列の次元を表現すると，何次元でも拡張することができるので今回はそのようにしました。}。宣言の時に，上で取り込んだ変数\verb|Lv|や\verb|N|を使うことができます。最大\verb|Lv|個の群，各群\verb|N|人のデータがあるという前提です。
	\item[parametersブロック] 全体平均\verb|gm|と，差分，誤差のSDをパラメータとしています。差分はここでは\verb|raw_delta|とし，上で宣言した\verb|Lv|をつかって\verb|Lv-1|個の要素を持つ配列にしています。\verb|raw_delta|って変な名前，と思うかもしれませんが，この後これを変換して$\delta$を作りますので，作る前の生(raw)の$\delta$という名前にしてみました\footnote{変数名は任意ですので，著者の命名法がわかりにくいなあと思ったら，自分で好きな変数名にしてもらって構いませんよ。その場合は，以後すべての変数を読み替えてくださいね。}。
	\item[transformed parametersブロック] 以後の話をわかりやすくするために，各群の平均\verb|mu|を構成することにします。\verb|mu|の数は水準数と同じ\verb|Lv|です。そして第i群の平均\verb|mu[i]|を作るために，\verb|mu[i]=gm+delta[i]|という書き方をしたいので，水準数と同じ数だけ\verb|delta[i]|を作りたいと思います。もとの\verb|raw_delta|が1,2,...,Lv-1個ありますので，\verb|delta|もLv-1番目まではそれと同じものをコピーします(\verb|for|文で繰り返し代入していってます。)。そして最後にLv番目の要素を，0からこれまでの要素をすべて足した(\verb|sum(raw_delta)|)ものを引くことで作っています。ちなみに\verb|sum|はstanの持っている総和の関数で，配列要素のすべてを足し合わせるというものです。こうして\verb|delta|がLv個できあがりましたので，各水準の平均値を\verb|mu[i]=gm+delta[i]|という数式で一般的に書くことができるようになりました。
	\item[modelブロック] 尤度のところに，群それぞれについて巡回する\verb|for|文，その中で1,2,3...N人まで巡回する\verb|for|文を組み込んで二重にぐるぐると回しています。事前分布は設定の通りです。
\end{description}
